The Temperature Control Laboratory (TC-Lab) provides transient heater and sensor
data for studying parameter estimation and experiment design. We compared
$n$-stage input-delay models in conductance--capacity and gain--time-constant
coordinates, separating formulation, initialization and numerical propagation.
The original four-parameter thermal model has an exact one-dimensional physical
parameter equivalence family under the modeled output and prescribed initial
conditions. The cold one-stage formulation represents its transfer behavior
with three parameters and a global pole-label ambiguity. Constant input with
equilibrated delay states leaves delay unobservable. Coordinate changes alter
conditioning without adding information. Cold joint calibration favors one
stage, while initialization changes the order ranking and transfer errors on
held-out records exceed calibration errors. Profiles, resampling and controlled
coverage comparisons show strong sensitivity to serial-error assumptions;
HAC intervals improve but do not restore nominal coverage in the longer
correlation scenario. Ambient/history and command-hold sensitivity further
limit the interpretation of narrow conditional intervals. Constrained local
input designs improve their declared information and fixed-mean separation
criteria. Refitting competitors removes equilibrated constant-input separation
and qualifies fixed-mean discrimination claims. The results support
identifiable transfer functions and explicit experimental assumptions as the
basis for inference. This draft and all scientific choices await human review.
